\SGARefTarget{sga3:xviii:section:3}\section*{3. Construction of a
Group from a Rational Law}
\addcontentsline{toc}{section}{3. Construction of a Group from a
Rational Law}

\SGARefTarget{sga3:xviii:paragraph:3.0}\paragraph{\SGARefLink{sga3:xviii:paragraph:3.0}{3.0}.}
Suppose given an \(S\)-prescheme \(X/S\) and a rational map
\[
  X\times_S X\dashrightarrow X
\]
over \(S\). We seek an \(S\)-group \(G/S\) and a birational map
relative to \(S\)\footnote{Editor's note: The notion of a
birational map (relative to \(S\)) should be made explicit, perhaps
in an addition to 1.5.2.}
\[
  X\dashrightarrow G
\]
that commutes with the composition laws. We treat only the case in
which \(X/S\) satisfies the following hypothesis:
\[
\tag{\(\diamond\)}\label{sga3:xviii:section:3:equation:diamond}
\begin{gathered}
  X/S\text{ is faithfully flat and of finite presentation,}\\
  \text{with separated fibers and without embedded components.}
\end{gathered}
\]
(The last two conditions hold for a group
prescheme.\footnote{Editor's note: Cf.\ VI\textsubscript{A}; the
reference is to be specified.})

We shall often omit the symbol \(S\) from fiber products.

Let \(X/S\) be a prescheme having the properties
\(\textup{(\(\diamond\))}\) above, and let \(W\) be a subprescheme
of finite presentation of \(X\times X\times X\) having the
following property:
\[
\tag{*}\label{sga3:xviii:section:3:equation:star}
\begin{gathered}
  \text{the three morphisms \(W\to X\times X\) induced by the
  projections}\\
  \text{from \(X^3\) onto \(X^2\) are relatively schematically
  dense open immersions over \(S\).}
\end{gathered}
\]

\paragraph{Notation.}
We shall use the following terminology. Given sections
\(a,b,c\in X(S')\), where \(S'\to S\) is arbitrary, such that
\((a,b,c)\in W(S')\), we write
\[
  c=ab,\qquad b=a^{-1}c,\qquad a=cb^{-1}.
\]

\SGARefTarget{sga3:xviii:def:3.0.1}\paragraph{\SGARefLink{sga3:xviii:def:3.0.1}{Definition 3.0.1}.}
\footnote{Editor's note: The numbering \SGARefLink{sga3:xviii:def:3.0.1}{3.0.1} and \SGARefLink{sga3:xviii:remark:3.0.2}{3.0.2} has been
introduced.}
Given a section \((a,b)\in X^2(S')\), we say that \(ab\) is
\emph{defined} if, and only if, there exists a section
\(c\in X(S')\) such that \((a,b,c)\in W(S')\); equivalently, if and
only if \((a,b)\in\operatorname{pr}_{12}W(S')\). Similarly, to say
that \(a^{-1}b\) or \(ab^{-1}\) is defined has the analogous
meaning, and we extend this terminology also to products of several
factors.

\SGARefTarget{sga3:xviii:remark:3.0.2}\paragraph{\SGARefLink{sga3:xviii:remark:3.0.2}{Remark 3.0.2}.}
Let us note at once the following fact. By \textup{(*)}, \(W\)
defines a rational map \(X^2\dashrightarrow X\) over \(S\), namely
the one given by \((a,b)\mapsto ab\). This rational map may well
have a domain of definition larger than
\(\operatorname{pr}_{12}W\). Nevertheless, we say that \(ab\) is
defined only if \((a,b)\in\operatorname{pr}_{12}W(S')\).

\SGARefTarget{sga3:xviii:def:3.1}\paragraph{\SGARefLink{sga3:xviii:def:3.1}{Definition 3.1}.}
A \emph{group germ} is an \(S\)-prescheme \(X/S\) having the
properties \(\textup{(\(\diamond\))}\) above, together with a
subprescheme \(W\) of finite presentation of
\(X\times X\times X\), having the following properties:
\begin{enumerate}
\item[\SGAInlineRefTarget{sga3:xviii:def:3.1:item:i}(i)] \(W\) satisfies property \textup{(*)} above.

\item[\SGAInlineRefTarget{sga3:xviii:def:3.1:item:ii}(ii)] For every section \(a\in X(S)\), the sets
\[
\begin{array}{ll}
  \operatorname{pr}_i\bigl((a\times X\times X)\cap W\bigr),
    & i=2,3,\\
  \operatorname{pr}_i\bigl((X\times a\times X)\cap W\bigr),
    & i=3,1,\\
  \operatorname{pr}_i\bigl((X\times X\times a)\cap W\bigr),
    & i=1,2,
\end{array}
\]
are schematically dense in \(X\) relative to \(S\), and this
assertion remains true after every base change \(S'\to S\).
(Condition \SGARefLink{sga3:xviii:def:3.1:item:i}{(i)} implies that these sets are open in \(X\).)
Intuitively, this says that \(ax\) is defined for sufficiently
general \(x\), and similarly for the other products.

\item[\SGAInlineRefTarget{sga3:xviii:def:3.1:item:iii}(iii)] The law is associative: if
\((a,b,c)\in X^3(S')\), for an arbitrary \(S'\to S\), and both
\((ab)c\) and \(a(bc)\) are defined, then
\[
  (ab)c=a(bc).
\]
\end{enumerate}

\paragraph{Remarks.}
\textup{(a)} In \SGARefLink{sga3:xviii:def:3.1:item:i}{(i)}, the condition that \(W\) be schematically
dense in \(X^2\) relative to \(S\) may be replaced by the condition
that, for every \(s\in S\), the fiber \(W_s\) be topologically
dense in \(X_s^2\), by the hypotheses on \(X\) and by \SGARefLink{sga3:xviii:prop:1.2}{1.2}.

\textup{(b)} Condition \SGARefLink{sga3:xviii:def:3.1:item:iii}{(iii)} is equivalent to the following.
Let \(V_1\) (respectively \(V_2\)) be the open subset of \(X^3\)
on which the rational map
\[
  (a,b,c)\longmapsto a(bc)
  \qquad
  \bigl(\text{respectively }(a,b,c)\longmapsto (ab)c\bigr)
\]
is defined\begingroup
\renewcommand{\thefootnote}{*}\footnote{In the sense explained in
\SGARefLink{sga3:xviii:def:3.0.1}{3.0.1}; one could also replace these open subsets by the domains of
definition (cf.\ \SGARefLink{sga3:xviii:def:1.5}{1.5}) of the rational maps under consideration.}
\addtocounter{footnote}{-1}
\endgroup. Then there exists an open subset
\(V\subseteq V_1\cap V_2\), schematically dense in \(X^3\)
relative to \(S\), on which the two maps agree. This follows from
\SGARefLink{sga3:xviii:prop:1.4}{1.4}, since \(V_1\) and \(V_2\) are schematically dense in \(X^3\)
relative to \(S\).

\textup{(c)} Hypothesis \SGARefLink{sga3:xviii:prop:1.4:item:ii}{(ii)} will be used below to ensure that
\(X\) is a subobject of the group prescheme that it defines. In
many cases one can deduce \SGARefLink{sga3:xviii:def:3.1:item:ii}{(ii)} from \SGARefLink{sga3:xviii:def:3.1:item:i}{(i)}, provided that \(X\) is
replaced by a relatively schematically dense open subset \(X'\)
and \(W\) by a relatively schematically dense open subset of
\(W\cap {X'}^3\). More precisely:

\SGARefTarget{sga3:xviii:prop:3.2}\paragraph{\SGARefLink{sga3:xviii:prop:3.2}{Proposition 3.2}.}
\textit{Suppose that every geometric fiber of \(X/S\) is
irreducible, and let \(W\) be a subprescheme of \(X^3\) satisfying
condition \SGARefLink{sga3:xviii:def:3.1:item:i}{(i)} of \SGARefLink{sga3:xviii:def:3.1}{3.1}. Then there exist a relatively schematically
dense open subset \(X'\subseteq X\) and a relatively schematically
dense open subset
\(W'\subseteq W\cap {X'}^3\)\begingroup
\renewcommand{\thefootnote}{\ensuremath{\dagger}}\footnote{The
French source prints \(W\cap(X'\times X')\) in this statement. The
proof immediately defines \(W'=W\cap {X'}^3\), which is the
dimensionally coherent reading used here.}
\addtocounter{footnote}{-1}
\endgroup such that the pair \((X',W')\) satisfies conditions
\SGARefLink{sga3:xviii:def:3.1:item:i}{(i)} and \SGARefLink{sga3:xviii:def:3.1:item:ii}{(ii)}. If \SGARefLink{sga3:xviii:def:3.1:item:iii}{(iii)} holds for \((X,W)\), then it holds for
\((X',W')\).}

\begin{proof}
Set
\[
  X'=\bigcap_{i=1}^{3}\operatorname{pr}_i W.
\]
Each \(\operatorname{pr}_iW\) is open in \(X\), because
\(W\to X^2\) is an open immersion and the projections
\(X^2\to X\) are flat and of finite presentation. Moreover,
\(\operatorname{pr}_iW\) is schematically dense in \(X\) relative
to \(S\), because \(W\) is so in \(X^2\) and the projections
\(X^2\to X\) are surjective; it is enough here to check
topological density. Put
\[
  W'=W\cap {X'}^3.
\]

To verify \SGARefLink{sga3:xviii:def:3.1:item:i}{(i)}, note that
\[
\begin{aligned}
  W'={}&
  \bigl(W\cap(X'\times X'\times X)\bigr)
  \cap\bigl(W\cap(X\times X'\times X')\bigr)\\
  &\cap\bigl(W\cap(X'\times X\times X')\bigr).
\end{aligned}
\]
Now
\[
  W\cap(X'\times X'\times X)
  \simeq \operatorname{pr}_{12}W\cap(X'\times X'),
\]
so this term is schematically dense in \(W\) relative to \(S\).
The other two terms on the right are likewise schematically dense
in \(W\), and hence \(W'\) is schematically dense in \(W\).
Consequently \(\operatorname{pr}_{ij}W'\) is schematically dense
in \(\operatorname{pr}_{ij}W\), hence in \(X\times X\), and
therefore in \(X'\times X'\).

To verify \SGARefLink{sga3:xviii:def:3.1:item:ii}{(ii)}, let \(a\in X'(S')\). After this base change, rename
\(S'\) as \(S\). We must show, for example, that
\[
  \operatorname{pr}_2\bigl((a\times X'\times X')\cap W'\bigr)
\]
is schematically dense in \(X'\) relative to \(S\). By \SGARefLink{sga3:xviii:cor:1.3}{1.3} it is
enough to check this fiber by fiber. Thus it is enough to treat the
case in which \(S\) is the spectrum of a field, and in that case it
is enough to show that the open subset is nonempty, since the
fibers of \(X/S\) are irreducible and have no embedded components.
Since
\[
  \operatorname{pr}_2\bigl((a\times X\times X)\cap W\bigr)
\]
is nonempty, because \(a\) is a section of \(X'\), this open subset
is dense in \(X\). We have
\[
  \operatorname{pr}_2\bigl((a\times X\times X)\cap W\bigr)
  =
  \operatorname{pr}_2\bigl((a\times X)
    \cap\operatorname{pr}_{12}W\bigr).
\]
Therefore
\[
\begin{aligned}
  &\operatorname{pr}_2\bigl((a\times X)
    \cap\operatorname{pr}_{12}W\bigr)\cap X'\\
  &\qquad =
  \operatorname{pr}_2\bigl((a\times X')
    \cap\operatorname{pr}_{12}W\bigr)
  =
  \operatorname{pr}_2\bigl((a\times X'\times X)\cap W\bigr)
\end{aligned}
\]
is dense in \(X\), and hence in \(X'\).

Similarly,
\(\operatorname{pr}_3((a\times X\times X')\cap W)\) is dense in
\(X\), hence dense in
\(\operatorname{pr}_3((a\times X\times X)\cap W)\). Equivalently,
\((a\times X\times X')\cap W\) is dense in
\((a\times X\times X)\cap W\), and therefore
\[
  \operatorname{pr}_2\bigl((a\times X\times X')\cap W\bigr)
\]
is dense in
\(\operatorname{pr}_2((a\times X\times X)\cap W)\), hence in
\(X\), and thus in \(X'\).

Finally,
\[
\begin{aligned}
  \operatorname{pr}_2\bigl((a\times X'\times X')\cap W'\bigr)
  &=
  \operatorname{pr}_2\bigl((a\times X'\times X')\cap W\bigr)\\
  &=
  \operatorname{pr}_2\bigl((a\times X'\times X)\cap W\bigr)
  \cap
  \operatorname{pr}_2\bigl((a\times X\times X')\cap W\bigr),
\end{aligned}
\]
which is indeed dense in \(X'\). This is what was required. The
other assertions in \SGARefLink{sga3:xviii:def:3.1:item:ii}{(ii)} follow by symmetry, and preservation of
condition \SGARefLink{sga3:xviii:def:3.1:item:iii}{(iii)} is immediate.\footnote{Editor's note: The
following sentence has been added.} This proves \SGARefLink{sga3:xviii:prop:3.2}{Proposition 3.2}.
\end{proof}

\SGARefTarget{sga3:xviii:paragraph:3.2.1}\paragraph{\SGARefLink{sga3:xviii:paragraph:3.2.1}{3.2.1}.}
\footnote{Editor's note: The numbering \SGARefLink{sga3:xviii:paragraph:3.2.1}{3.2.1}, as well as \SGARefLink{sga3:xviii:paragraph:3.2.2}{3.2.2},
has been introduced.}
Fix a group germ \((X,W)\) over \(S\). We first make some
preliminary observations about the situation, collected below as
rules. These rules will often be used without explicit mention in
what follows.

Let \(a\in X(S')\), and after base change rename \(S'\) as \(S\).
There is then a birational map \(\varphi\) over \(S\) from \(X\)
to itself, sending a section \(x\) to \(ax\) whenever the latter is
defined. By \SGARefLink{sga3:xviii:def:3.1:item:ii}{3.1(ii)}, the domain of definition of \(\varphi\)
contains the relatively schematically dense open subset
\[
  \operatorname{pr}_2\bigl((a\times X\times X)\cap W\bigr),
\]
and \(\varphi\) defines an isomorphism from this open subset onto
the open subset
\[
  \operatorname{pr}_3\bigl((a\times X\times X)\cap W\bigr),
\]
on which \(a^{-1}x\) is defined. This observation is generalized
in the evident way in Rule 1 below.

\SGARefTarget{sga3:xviii:rule:1}\paragraph{Rule 1.}
Let \(P=P(x,t_1,\ldots,t_n)\) be a ``product'' of the symbols
\(x,t_1,\ldots,t_n\), obtained recursively as follows:
\(P_0=x\), and \(P_{i+1}\) is one of the expressions
\[
  P_i t,\qquad tP_i,\qquad P_i^{-1}t,\qquad
  tP_i^{-1},\qquad P_i t^{-1},\qquad t^{-1}P_i,
\]
where \(t\) is one of the \(t_j\); and \(P_r=P\). Let
\(a_1,\ldots,a_n\in X(S')\). Then there exists a relatively
schematically dense open subset \(U\) of \(X_{S'}\) such that the
product \(P(x,a_1,\ldots,a_n)\) is defined (in the sense of
\SGARefLink{sga3:xviii:remark:3.0.2}{Remark 3.0.2}) for a section \(x\in X(S'')\) if and only if
\(x\in U(S'')\), and the map
\[
  x\longmapsto P(x,(a))
\]
gives an isomorphism from \(U\) onto another relatively
schematically dense open subset of \(X_{S'}\), denoted
\(P(U,(a))\).

\SGARefTarget{sga3:xviii:rule:2}\paragraph{Rule 2.}
Let \(a,b\in X(S')\). Then:
\begin{align*}
  &\text{if \(ab\) is defined, so is \(a^{-1}(ab)\), and }
    a^{-1}(ab)=b;\\
  &\text{if \(a^{-1}b\) is defined, so is \(a(a^{-1}b)\), and }
    a(a^{-1}b)=b;\\
  &\text{if \(ba^{-1}\) is defined, so is \((ba^{-1})a\), and }
    (ba^{-1})a=b.
\end{align*}

\SGARefTarget{sga3:xviii:rule:3}\paragraph{Rule 3.}
Let \(a,b,b'\in X(S')\). If
\[
  ab=ab',\qquad ba=b'a,\qquad a^{-1}b=a^{-1}b',
  \qquad\text{or}\qquad ba^{-1}=b'a^{-1},
\]
then \(b=b'\). It is understood here that an equality entails that
both of its sides are defined.

\SGARefTarget{sga3:xviii:rule:4}\paragraph{Rule 4.}
Let \(a,b,c\in X(S')\). Whenever both sides are defined, one has
\[
  a\bigl((ba)^{-1}c\bigr)=b^{-1}c.
\]
Similarly,
\[
  \bigl(c(ab)^{-1}\bigr)a=cb^{-1},\qquad
  a^{-1}\bigl((ab^{-1})c\bigr)=b^{-1}c,\qquad
  \bigl(c(b^{-1}a)\bigr)a^{-1}=cb^{-1}.
\]

\SGARefTarget{sga3:xviii:rule:5}\paragraph{Rule 5.}
All the following associativity laws hold whenever both sides are
defined:
\[
\begin{aligned}
  (a^{-1}b)c&=a^{-1}(bc),&
  (ab^{-1})c&=a(b^{-1}c),&
  (ab)c^{-1}&=a(bc^{-1}),\\
  (ab)^{-1}c&=b^{-1}(a^{-1}c),&
  (a^{-1}b)c^{-1}&=a^{-1}(bc^{-1}),&
  (ab^{-1})c^{-1}&=a(cb)^{-1}.
\end{aligned}
\]

\SGARefTarget{sga3:xviii:paragraph:3.2.2}\paragraph{\SGARefLink{sga3:xviii:paragraph:3.2.2}{3.2.2}.
Verification of the rules.}
\textup{\SGARefLink{sga3:xviii:section:1}{(1)}} This follows by an evident induction on the length
\(r\) of \(P\), the case \(r=1\) being a direct consequence of
\SGARefLink{sga3:xviii:def:3.1:item:ii}{3.1(ii)}.

\textup{\SGARefLink{sga3:xviii:section:2}{(2)}} This is immediate from the definition.

\textup{\SGARefLink{sga3:xviii:section:3}{(3)}} Indeed, by Rule 2, in the first case for example,
\[
  b=a^{-1}(ab)=a^{-1}(ab')=b'.
\]

\textup{(4)} Let us verify the first relation, for example. On the
right-hand side, left multiplication by \(b\) is defined and gives
\(c\), by Rule 2. Suppose that it is also defined on the
left-hand side. Then
\[
  b\bigl(a((ba)^{-1}c)\bigr)
  =(ba)\bigl((ba)^{-1}c\bigr)
  =c.
\]
Indeed, \(ba\) is defined by hypothesis, since it occurs in the
expression. The middle term is therefore defined and equal to
\(c\) by Rule 2, and it is equal to the left-hand term by
associativity (cf.\ \SGARefLink{sga3:xviii:def:3.1:item:iii}{3.1(iii)}). Rule 3 thus implies that the desired
equality holds whenever this multiplication by \(b\) is defined.

Now fix \(a\) and \(b\). Rule 1 implies that
\(b(a((ba)^{-1}c))\) is defined for \(c\) ``in'' a relatively
schematically dense open subset \(U\) of \(X\). Hence, on this
relatively schematically dense open subset, the two rational maps
\[
  c\longmapsto b^{-1}c
  \qquad\text{and}\qquad
  c\longmapsto a((ba)^{-1}c)
\]
are equal. By \SGARefLink{sga3:xviii:prop:1.4}{1.4}, they are equal on every common domain of
definition, which proves the desired result.

\textup{(5)} The argument is of the same kind as the preceding
one. For example, let us verify
\[
  (ab^{-1})c^{-1}=a(cb)^{-1}.
\]
If right multiplication by \(c\) is defined on the right-hand
side, equality follows from Rule 4. Since it is enough to verify
such a formula on a relatively schematically dense open subset, we
are reduced to the case in which this multiplication is defined.

\SGARefTarget{sga3:xviii:paragraph:3.2.3}\paragraph{\SGARefLink{sga3:xviii:paragraph:3.2.3}{3.2.3}.}
Consider the relation \(R\) on \(X\times X\) defined as follows.
For \(a,b,a',b'\in X(S')\), one has
\[
  (a,b)\sim(a',b')\pmod{R(S')}
\]
if and only if there exist an fppf covering \(S''\to S'\) and a
section \(x\in X(S'')\) such that \((xa)b\) and \((xa')b'\) are
defined and equal. Then \(R\) is an equivalence relation.

Indeed, this relation is evidently symmetric. By Rule 1, the
product \((xa)b\) is defined when \(x\) is ``in'' a suitable
relatively schematically dense open subset. Thus \SGARefLink{sga3:xviii:prop:1.7}{1.7} asserts that
there exist an fppf covering \(S''\to S'\) and an
\(x\in X(S'')\) for which \((xa)b\) is defined. The relation is
therefore reflexive. Its transitivity follows from the next lemma.

\SGARefTarget{sga3:xviii:lemma:3.3}\paragraph{\SGARefLink{sga3:xviii:lemma:3.3}{Lemma 3.3}.}
\textit{Let \(x,y,a,b,a',b'\) be sections of \(X(S')\) such that
\((xa)b\), \((xa')b'\), \((ya)b\), and \((ya')b'\) are defined. If
\[
  (xa)b=(xa')b',
\]
then
\[
  (ya)b=(ya')b'.
\]}

Indeed, the lemma says that one may test
\((a,b)\sim(a',b')\) using any \(x\in X(S')\), with
\(S'\to S\) an fppf covering, for which both products are defined.
Given
\[
  a,b,a',b',a'',b''\in X(S),
\]
Rules 1, \SGARefLink{sga3:xviii:prop:1.1}{1.1}, and \SGARefLink{sga3:xviii:prop:1.7}{1.7} allow us to find an fppf covering
\(S'\to S\) and a section \(x\in X(S')\) such that the three
products in question are defined; transitivity follows.

\begin{proof}[Proof of \SGARefLink{sga3:xviii:lemma:3.3}{Lemma 3.3}]
Let us begin by writing formally
\[
\begin{aligned}
  (za)b
  &=\bigl(((zx^{-1})x)a\bigr)b
   =\bigl((zx^{-1})(xa)\bigr)b
   =(zx^{-1})\bigl((xa)b\bigr),\\
  (za')b'
  &=\bigl(((zx^{-1})x)a'\bigr)b'
   =\bigl((zx^{-1})(xa')\bigr)b'
   =(zx^{-1})\bigl((xa')b'\bigr).
\end{aligned}
\]
The appropriate rules show that these equalities do hold whenever
their terms are defined.\footnote{Editor's note: To be made
precise\ldots} It follows that
\[
  (za)b=(za')b'
\]
whenever all these expressions are defined. Moreover, by Rule 1
and the hypotheses already made, these expressions are defined
when \(z\in X(S'')\) belongs to \(V(S'')\), where \(V\) is a
certain relatively schematically dense open subset of \(X\) (we
have taken \(S'=S\)). Therefore the two rational self-maps of
\(X\) given by
\[
  z\longmapsto(za)b
  \qquad\text{and}\qquad
  z\longmapsto(za')b'
\]
are equal, and consequently
\[
  (ya)b=(ya')b'.
\]
\end{proof}

\SGARefTarget{sga3:xviii:lemma:3.4}\paragraph{\SGARefLink{sga3:xviii:lemma:3.4}{Lemma 3.4}.}
\textit{Consider the rational map
\[
  \varphi\colon X^3\dashrightarrow X
\]
over \(S\), defined by
\[
  (a,b,c)\longmapsto c^{-1}(ab).
\]
Let \(U\) be the domain of definition of \(\varphi\), and let
\(\Gamma\) be the graph of the morphism \(f\colon U\to X\)
induced by \(\varphi\), regarded as a subprescheme of \(X^4\).
Then a section \((a,b,c,d)\in X^4(S)\) belongs to \(\Gamma(S)\)
if and only if
\[
  (a,b)\sim(c,d).
\]}

\begin{proof}
First observe that the rational map \(\varphi\) is also given by
the formula
\[
  (a,b,c)\longmapsto(xc)^{-1}\bigl((xa)b\bigr)
\]
for an arbitrary section \(x\in X(S)\). Equivalently,
\[
  c^{-1}(ab)=(xc)^{-1}\bigl((xa)b\bigr)
\]
whenever both sides are defined. We leave the verification to the
reader.

We now show that \(\varphi\) is defined at a section
\((a,b,c)\in X^3(S)\) if and only if there exists a \(d\) such
that \((a,b)\sim(c,d)\), and that in this case
\[
  d=\varphi(a,b,c).
\]
Suppose first that \((a,b)\sim(c,d)\). To verify that
\(\varphi\) is defined, we may make an fppf-covering base change
(\SGARefLink{sga3:xviii:prop:1.6}{Proposition 1.6}), and may therefore suppose that there exists a
section \(x\in X(S')\) such that \((xa)b\) and \((xc)d\) are
defined and equal. It follows from Rule 2 that
\[
  (xc)^{-1}\bigl((xa)b\bigr)
\]
is defined and equal to \(d\).

Conversely, suppose that \(\varphi\) is defined at
\((a,b,c)\), and put \(d=\varphi(a,b,c)\). Choose an
fppf covering \(S'\to S\) and a section \(x\in X(S')\) such that
\((xa)b\) and \((xc)d\) are defined. We wish to show that they
are equal. It is enough to show that the two rational self-maps of
\(X\) over \(S\), given by
\[
  b\longmapsto(xa)b
  \qquad\text{and}\qquad
  b\longmapsto(xc)\varphi(a,b,c),
\]
are equal; this follows from the observation in the first
paragraph.
\end{proof}

\SGARefTarget{sga3:xviii:prop:3.5}\paragraph{\SGARefLink{sga3:xviii:prop:3.5}{Proposition 3.5}.}
\textit{The equivalence relation
\[
  R\rightrightarrows X^2
\]
is representable, and it is a flat relation of finite
presentation; that is, the projections of \(R\) onto \(X^2\) are
flat morphisms of finite presentation.}

\begin{proof}
We may suppose that \(S\) is affine. Since \((X,W)\) is of finite
presentation over \(S\), we may descend the entire situation to an
\(S\) of finite type over \(\operatorname{Spec}\mathbb Z\), hence
to a noetherian \(S\). We may therefore suppose that \(S\) is
noetherian. It is then immediate that the graph \(\Gamma\) of
\SGARefLink{sga3:xviii:lemma:3.4}{3.4} is of finite presentation over \(X^2\). The projection
\(\left.\operatorname{pr}_{12}\right|_\Gamma\) is flat because
\[
  \Gamma\xrightarrow{\sim}\operatorname{pr}_{123}\Gamma,
\]
the latter being an open subset of \(X^3\), and because the
projection
\[
  \operatorname{pr}_{12}\colon X^3\longrightarrow X^2
\]
is flat, since \(X\) is flat over \(S\).

I claim that \(\Gamma\) represents \(R\). There is something to
prove here, because the domain of definition of a rational map
does not in general commute with base change. Let \(S''\to S\).
What is clear from \SGARefLink{sga3:xviii:lemma:3.4}{3.4} applied over \(S''\) is that
\[
  R(S'')\supseteq\Gamma(S''),
\]
because \(\varphi\times_S S''\) is certainly defined on
\(U\times_S S''\). Now let
\[
  (a,b,c,d)\in R(S'').
\]
We must prove that \((a,b,c,d)\in\Gamma(S'')\). This can be
checked locally, say for the étale topology. We may therefore
suppose that \(S''\) is strictly local, that is, the spectrum of a
henselian ring with separably closed residue field. Moreover, by
applying \SGARefLink{sga3:xviii:prop:1.6}{1.6} and the usual passage-to-the-limit arguments, we are
reduced to the case in which \(S\) is strictly local and
\(S''\to S\) is local.

Suppose that there is a section \(x\in X(S)\) such that over
\(S''\) the products \((xa)b\) and \((xc)d\) are defined. This
implies that
\[
  (xc)^{-1}\bigl((xa)b\bigr)
\]
is defined and equal to \(d\). Now there is a relatively
schematically dense open subset \(V\) of \(X^3\) such that
\((xc)^{-1}((xa)b)\) is defined if and only if
\((a,b,c)\in V(S'')\), and \(V\subseteq U\). Consequently,
\((a,b,c,d)\in\Gamma(S'')\) if such an \(x\) exists. By \SGARefLink{sga3:xviii:prop:1.6}{1.6}, we
may make an fppf-covering base change \(S'\to S\) in order to find
such an \(x\). Since \(S\) is strictly local, we can
find\begingroup
\renewcommand{\thefootnote}{*}\footnote{Combining Exposé
VI\textsubscript{A}, 1.1.1 and EGA IV\textsubscript{4}, 17.16.2.}
\addtocounter{footnote}{-1}
\endgroup{} a faithfully flat, local, finite morphism
\(S'\to S\) and a section \(x\in X(S')\) that passes through an
arbitrary closed point of the closed fiber of \(X/S\). To say
that \((xa)b\) and \((xc)d\) are defined means that, over
\(S''\), \(x\) lies in a certain relatively schematically dense
open subset; this can be checked on the closed fiber of
\(X_{S''}\). Thus the required section exists.
\end{proof}

Let \(G\) now be the quotient of \(X^2\) by \(R\), as a sheaf for
the fppf topology. We define a composition law on \(G\) as
follows. Let
\[
  (g,g')\in G(S')
\]
be represented by a section
\[
  ((a,b),(c,d))\in (X^2\times_S X^2)(S''),
\]
where \(S''\to S'\) is an fppf covering. Suppose moreover that
there is a section \(x\in X(S'')\) such that \(a(b(cx))\) and
\(x^{-1}d\) are defined; this is permitted by Rule 1 and \SGARefLink{sga3:xviii:prop:1.7}{1.7}. We
define \(gg'\) to be the class in \(G(S')\) represented by the
section
\[
  \bigl(a(b(cx)),x^{-1}d\bigr)\in X^2(S'').
\]

Let us verify that \(gg'\) is independent of the choice of the
section \(x\) and of the representative
\(((a,b),(c,d))\). Let
\[
  (a',b')\sim(a,b),\qquad (c',d')\sim(c,d),
\]
and let \(x'\) be such that \(a'(b'(c'x'))\) and
\({x'}^{-1}d'\) are defined. We may suppose that all these are
sections over \(S''\). We must prove that
\[
  \bigl(a(b(cx)),x^{-1}d\bigr)
  \sim
  \bigl(a'(b'(c'x')),{x'}^{-1}d'\bigr);
\]
that is, for a suitable fppf covering \(S'''\to S''\) and a
suitable section \(z\in X(S''')\), we must have
\[
 \bigl(z(a(b(cx)))\bigr)(x^{-1}d)
 =
 \bigl(z(a'(b'(c'x')))\bigr)({x'}^{-1}d').
\]
Whenever all the products involved are defined,
\[
\begin{aligned}
 \bigl(z(a(b(cx)))\bigr)(x^{-1}d)
 &=\bigl((za)(b(cx))\bigr)(x^{-1}d)\\
 &=\bigl(((za)b)(cx)\bigr)(x^{-1}d)\\
 &=\bigl((((za)b)c)x\bigr)(x^{-1}d)\\
 &=\bigl(((za)b)c\bigr)\bigl(x(x^{-1}d)\bigr)\\
 &=\bigl(((za)b)c\bigr)d,
\end{aligned}
\]
and the same identities hold with primes. By Rule 1 and \SGARefLink{sga3:xviii:prop:1.7}{1.7},
such a \(z\) exists. It therefore remains to prove that
\[
  \bigl(((za)b)c\bigr)d
  =
  \bigl(((za')b')c'\bigr)d'.
\]
But \((za)b=(za')b'\), since
\((a,b)\sim(a',b')\) by \SGARefLink{sga3:xviii:lemma:3.3}{3.3}; the desired equality now follows
from \((c,d)\sim(c',d')\).

Consider the natural morphism
\[
  i\colon X\longrightarrow G
\]
defined as follows. For \(a\in X(S')\), choose an fppf covering
\(S''\to S'\) and a section \(b\in X(S'')\) such that
\(ab^{-1}\) is defined, and put
\[
  i(a)=\text{the class in \(G\) of }\bigl(ab^{-1},b\bigr).
\]
It is readily verified that this class, which a priori lies in
\(G(S'')\), is independent of the choice of \(b\), and therefore
defines a well-determined element of \(G(S')\).

The reader may enjoy verifying the following proposition.

\SGARefTarget{sga3:xviii:prop:3.6}\paragraph{\SGARefLink{sga3:xviii:prop:3.6}{Proposition 3.6}.}
\textit{The morphism \(i\) commutes with the composition laws on
\(X\) and \(G\). In other words, if \(a,b\in X(S')\) are
sections for which \(ab\) is defined, then
\[
  i(a)i(b)=i(ab).
\]}

The goal of this section is the following theorem.

\SGARefTarget{sga3:xviii:thm:3.7}\paragraph{\SGARefLink{sga3:xviii:thm:3.7}{Theorem 3.7}.}
\footnote{Editor's note: If \(X\) is smooth and separated over
\(S\), faithfully flat and of finite presentation over \(S\),
then \(G/S\) is representable by a smooth \(S\)-group scheme of
finite type over \(S\). This is Theorem 6.6.1 of
\emph{Néron Models}, by Bosch--Lütkebohmert--Raynaud, Springer
(1990).}
\textit{Let \((X,W)\) be a group germ over \(S\), with \(X/S\)
faithfully flat and of finite presentation, and with separated
fibers having no embedded component. With the notation above:}
\begin{enumerate}
\item[\SGAInlineRefTarget{sga3:xviii:thm:3.7:item:i}(i)] \textit{\(G\) is a sheaf of groups.}
\item[(ii)] \textit{\(i\colon X\to G\) is representable by an
open immersion.}
\item[\SGAInlineRefTarget{sga3:xviii:theorem:3-7:item:iii}\SGAInlineRefTarget{sga3:xviii:thm:3.7:item:iii}(iii)] \textit{\(G\) is representable fppf-locally over
\(S\).}
\item[\SGAInlineRefTarget{sga3:xviii:thm:3.7:item:iv}(iv)] \textit{If \(G/S\) is representable, then it is a
flat group of finite presentation, and
\(i\colon X\to G\) is schematically dense relative to \(S\).}
\end{enumerate}

Observe that \(G\) is plainly characterized by properties
\textup{\SGARefLink{sga3:xviii:thm:3.7:item:i}{(i)}}--\textup{\SGARefLink{sga3:xviii:thm:3.7:item:iv}{(iv)}}; one may therefore forget its
explicit construction.

The proof proceeds in several stages.

\SGARefTarget{sga3:xviii:lemma:3.8}\paragraph{\SGARefLink{sga3:xviii:lemma:3.8}{Lemma 3.8}.}
\begin{enumerate}
\item[\SGAInlineRefTarget{sga3:xviii:lemma:3.8:item:i}(i)] \textit{Let \(c\) be a section of \(X(S)\). Then the
morphism from the prescheme \(c\times X\), which is
\(S\)-isomorphic to \(X\), into \(G\), given by
\[
  c\times X\lhook\joinrel\longrightarrow X\times X
  \longrightarrow G,
\]
is a monomorphism.}

\item[\SGAInlineRefTarget{sga3:xviii:lemma:3.8:item:ii}(ii)] \textit{Let \(\{c_i\}_{i\in I}\) be sections of
\(X(S)\), put
\[
  Z=\coprod_i c_i\times X,
\]
and let \(R'\rightrightarrows Z\) be the equivalence relation
induced on \(Z\) by the evident morphism \(Z\to X^2\). Then
\(R'\) is a ``gluing'' of the \(c_i\times X\simeq X\).}
\end{enumerate}

\begin{proof}
\textup{\SGARefLink{sga3:xviii:lemma:3.8:item:i}{(i)}} For two sections \((c,a)\) and \((c,a')\) of
\((c\times X)(S')\) to have the same image in \(G\), one must
have \((c,a)\sim(c,a')\); that is, for a suitable section
\(x\in X(S'')\), where \(S''\to S'\) is an fppf covering,
\[
  (xc)a=(xc)a'
\]
\begingroup
\renewcommand{\thefootnote}{\ensuremath{\dagger}}\footnote{The
French source prints \(\sim\) between these two sections of
\(X\). Equality is the relation-defining condition and is
required for the following application of Rule 3.}
\addtocounter{footnote}{-1}
\endgroup.
It follows from Rule 3 that \(a=a'\).

\textup{\SGARefLink{sga3:xviii:lemma:3.8:item:ii}{(ii)}} Let \(c_i,c_j\) be two sections, and consider the
birational self-map \(\psi_{ji}\) of \(X\) over \(S\), given by
\[
  x\longmapsto c_j^{-1}(c_i x).
\]
\begingroup
\renewcommand{\thefootnote}{\ensuremath{\ddagger}}\footnote{The
French source introduces this map as \(\psi_{ji}\), but calls it
\(\varphi_{ji}\) thereafter. We use \(\psi_{ji}\) consistently.}
\addtocounter{footnote}{-1}
\endgroup
For \(y\in X(S)\), this is the same map as
\[
  x\longmapsto (yc_j)^{-1}\bigl((yc_i)x\bigr),
\]
as is readily seen. Moreover, \(\psi_{ji}\) is defined at a
section \(b\in X(S')\) if and only if there exists a section
\(b'\) such that
\[
  (c_i,b)\sim(c_j,b'),
\]
and then \(b'=\psi_{ji}(b)\). Let \(U_{ji}\) be the domain of
definition of \(\psi_{ji}\) over \(S\). It remains to prove that
this domain of definition is universal: if \(b\in X(S'')\), for
an arbitrary \(S''\to S\), and \(\psi_{ji}\) is defined at
\(b\), then \(b\in U_{ji}(S'')\). Equivalently, one must prove
that if \(b,b'\in X(S'')\) satisfy
\[
  (c_i,b)\sim(c_j,b'),
\]
then \(b\in U_{ji}(S'')\). We leave to the reader the
verification of this fact, which is analogous to that in \SGARefLink{sga3:xviii:prop:3.5}{3.5}.
\end{proof}

\SGARefTarget{sga3:xviii:lemma:3.9}\paragraph{\SGARefLink{sga3:xviii:lemma:3.9}{Lemma 3.9}.}
\textit{Suppose that \(\{c_i\}_{i\in I}\) are sections of
\(X(S)\) such that
\[
  \coprod_i c_i\times X\longrightarrow G
\]
is surjective as a morphism of sheaves. Then \(G\) is
representable, flat, and of finite presentation over \(S\), and
the structural morphism
\[
  X^2\longrightarrow G
\]
is flat and of finite presentation.}

\begin{proof}
The representability of \(G\) is an immediate consequence of
\SGARefLink{sga3:xviii:lemma:3.8}{Lemma 3.8}, and it follows that
\[
  \coprod_i c_i\times X\longrightarrow G
\]
is an open covering. To prove that \(X^2\to G\) is flat, it is
enough to do so locally, hence to prove that the rational map
\[
  X^2\dashrightarrow c_i\times X
\]
induces a flat morphism on its domain of definition. This rational
map is given by
\[
  (a,b)\longmapsto
  \bigl(c_i,\bigl((xc_i)^{-1}(xa)\bigr)b\bigr),
\]
where \(x\) is an arbitrary section. If it is defined at
\((a,b)\in X^2(S')\), one can find an fppf covering
\(S''\to S'\) and a section \(x\in X(S'')\) such that
\(\bigl((xc_i)^{-1}(xa)\bigr)b\) is defined. It is then readily
seen that this morphism is flat. The same argument shows that it
is locally of finite presentation, and hence is fppf-covering.

By construction the relation \(R\) is effective. \SGARefLink{sga3:xviii:prop:3.5}{Proposition 3.5}
therefore shows that \(X^2\to G\) becomes a morphism of finite
presentation after the fppf-covering base change
\(G\leftarrow X^2\). Thus \(X^2\to G\) is of finite presentation.

Let us prove that \(G\to S\) is flat and of finite presentation.
It is flat and locally of finite presentation because \(G\) is
covered by the \(c_i\times X\simeq X\). Moreover, \(X^2/S\) is
quasi-compact and \(X^2\to G\) is surjective, so \(G/S\) is
quasi-compact. To prove that \(G\to S\) is quasi-separated, note
the following cartesian diagram:
\[
\begin{tikzcd}[column sep=large,row sep=large]
  R \arrow[r,"\alpha"] \arrow[d,"\gamma"'] &
  X^2\times_S X^2 \arrow[d,"\beta"]\\
  G \arrow[r,"\Delta"] & G\times_S G.
\end{tikzcd}

\]
Here \(\gamma\) is surjective and \(\alpha,\beta\) are
quasi-compact. Hence \(\beta\alpha\) is quasi-compact, and
therefore so is \(\Delta\).\footnote{Editor's note: Another way
to conclude here is by faithfully flat descent
\textup{(EGA IV\textsubscript{2}, 2.7.1)}, since \(\beta\) is an
fppf covering.}
\end{proof}

\SGARefTarget{sga3:xviii:lemma:3.10}\paragraph{\SGARefLink{sga3:xviii:lemma:3.10}{Lemma 3.10}.}
\textit{Let \(\{c_i\}_{i\in I}\) be sections of \(X(S)\). For
\[
  \coprod_i c_i\times X\longrightarrow G
\]
to be surjective as a morphism of fppf sheaves, it is enough that
the following condition hold: for every \(S'\to S\) and every
\((a,b)\in X^2(S')\), there exist an open covering
\(\{S'_\nu\}_{\nu\in N}\) of \(S'\) and a function
\[
  N\longrightarrow I,\qquad \nu\longmapsto i(\nu),
\]
such that \(\bigl(c_{i(\nu)}^{-1}a\bigr)b\) is defined over
\(S'_\nu\).}

\begin{proof}
Let \(S''\to S\) be arbitrary, and let \(g\in G(S'')\). Choose
an fppf covering \(S'\to S''\) and a section
\((a,b)\in X^2(S')\) representing \(g\). Take the open covering
\(\{S'_\nu\}\) supplied by the hypothesis. Over each \(S'_\nu\)
one has
\[
  (a,b)\sim
  \bigl(c_{i(\nu)},\bigl(c_{i(\nu)}^{-1}a\bigr)b\bigr).
\]
Thus, locally over \(S'_\nu\), the element \(g\) is represented
by a section of \(c_{i(\nu)}\times X\). Since the family
\(\{S'_\nu\to S''\}\) is fppf-covering, the desired surjectivity
follows.
\end{proof}

\SGARefTarget{sga3:xviii:lemma:3.11}\paragraph{\SGARefLink{sga3:xviii:lemma:3.11}{Lemma 3.11}.}
\textit{Let \(Y/S\) be a prescheme of finite presentation, and
let \(\{a_i\}_{i\in I}\) be sections of \(Y(S)\). Let
\(s_0,s_1\) be points of \(S\), with \(s_0\) a specialization of
\(s_1\), and let \(Y_j\) be the fiber of \(Y/S\) at \(s_j\).
Let \(C_j\) be the closure in \(Y_j\) of the set of points
\(\{a_i(s_j)\}_{i\in I}\). Then
\[
  \dim C_1\geq \dim C_0.
\]}

\begin{proof}
It is enough to verify the assertion after a base change
\(S'\to S\) together with chosen points \(s'_0,s'_1\) such that
\(s'_j\mapsto s_j\) and \(s'_0\) is a specialization of
\(s'_1\). By EGA II, 7.1.4, we are therefore reduced to the case
in which \(S\) is the spectrum of a valuation ring \(A\),
\(s_0\) is its closed point, and \(s_1\) is its generic point.

Let \(V\) be the closure of \(C_1\) in \(Y\). Clearly
\(C_0\subseteq V\). The lemma is consequently a special case of
the ``well-known'' fact that an irreducible closed subprescheme
\(V\) of a prescheme \(Y/S\) of finite presentation satisfies
\[
  \dim(V\times_S s_1)\geq \dim(V\times_S s_0)
\]
when \(S\) is the spectrum of a valuation ring and
\(V\times_S s_1\neq\varnothing\) (EGA IV, 13.1.6).
\end{proof}

\SGARefTarget{sga3:xviii:lemma:3.12}\paragraph{\SGARefLink{sga3:xviii:lemma:3.12}{Lemma 3.12}.}
\textit{Suppose that \(S\) is the spectrum of a local ring, with
closed point \(s_0\), and let \(\{c_i\}_{i\in I}\) be sections
such that the closure \(C_0\), in the closed fiber \(X_0\), of
the set \(\{c_i(s_0)\}_{i\in I}\) has dimension
\[
  \dim C_0=\dim X_0=n.
\]
Then the condition of \SGARefLink{sga3:xviii:lemma:3.9}{Lemma 3.9} is satisfied.}

\begin{proof}
First note that all the fibers of \(X/S\) have the same
dimension. This follows from EGA IV, 12.1.1(i), and from the fact
that \(X\) has a rational composition law. Lemma 3.11 therefore
implies that, for every morphism \(S_1\to S\), with \(S_1\) the
spectrum of a field, the closure of
\(\{c_i\times_S S_1\}_{i\in I}\) in \(X_{S_1}\) has dimension
\(n\).

Let us verify the condition of \SGARefLink{sga3:xviii:lemma:3.10}{Lemma 3.10}. Let
\((a,b)\in X^2(S')\). For \((c_i^{-1}a)b\) to be defined, it is
necessary and sufficient that \(c_i\) lie in a certain open
subset \(U\subseteq X_{S'}\) that is schematically dense relative
to \(S'\) (Rule 1). We must show that this holds for a suitable
\(i\), locally on \(S'\). It is therefore enough to treat the
case in which \(S'\) is the spectrum of a local ring; membership
of \(c_i\) in \(U(S')\) can then be checked on the closed fiber.
We are thus reduced to \(S'=\operatorname{Spec}k\), where \(k\)
is a field.

With the notation above, take \(S'=S_1\). Then
\[
  \dim C_1=\dim X_{S_1},
\]
and \(U\) is relatively schematically dense in \(X_{S_1}\).
Hence \(U\cap C_1\neq\varnothing\), and therefore
\[
  U\cap\{c_i\times_S S_1\}_{i\in I}\neq\varnothing.
\]
This proves the assertion.
\end{proof}

\noindent\textit{Proof of \SGARefLink{sga3:xviii:thm:3.7}{Theorem 3.7}.}
The proof is now easy. First note the following consequence of
the finiteness assertion in \SGARefLink{sga3:xviii:lemma:3.9}{Lemma 3.9}. Let \(\{A_i\}\) be a
directed system of rings over \(S\), put
\[
  \widetilde A=\varinjlim_i A_i,
\]
and suppose that the hypotheses of \SGARefLink{sga3:xviii:lemma:3.9}{Lemma 3.9} hold over
\(S=\operatorname{Spec}\widetilde A\). Then the object
representing the quotient \(G\) of
\[
  R\rightrightarrows X^2
\]
descends to some \(S_i=\operatorname{Spec}A_i\), together with
the finiteness and flatness properties stated in Lemma 3.9. This
is the usual passage to the limit (EGA IV, 8 and 11).

It follows that, in proving \SGARefLink{sga3:xviii:theorem:3-7:item:iii}{Theorem 3.7(iii)} and \SGARefLink{sga3:xviii:thm:3.7:item:iv}{(iv)}, we may
restrict ourselves to the case
\[
  S=\operatorname{Spec}A,
\]
where \(A\) is a strictly local ring. Let \(x_0\) be a closed
point of the closed fiber \(X_0\) of \(X/S\). There
exists\begingroup
\renewcommand{\thefootnote}{\ensuremath{\ast}}\footnote{Cf. the
note at the foot of page 15 in Exposé VII.}
\addtocounter{footnote}{-1}
\endgroup{} a local, finite, free extension \(A'\) of \(A\), and
a section of
\[
  X'=X\times_S S'
\]
passing through the unique point \(x'_0\) of \(X'\) lying over
\(x_0\). Notice that \(X'_0\to X_0\) is radicial because \(S\)
is strictly local, and hence the residue field of \(A\) is
separably closed.

It follows that there exists a directed system \(\{A_i\}\) of
local rings, flat and finite over \(A\), such that, on putting
\[
  \widetilde A=\varinjlim_i A_i,\qquad
  \widetilde X=X\times_S\operatorname{Spec}\widetilde A,
\]
the prescheme \(\widetilde X\) has a set of sections inducing a
dense set on its closed fiber \(\widetilde X_0\). Indeed, for
each closed point \(x_0\) of \(X_0\), choose an extension
\(A(x_0)\) such that the corresponding base change of \(X\)
admits a section ``passing through \(x_0\)'', and take for the
directed system the finite tensor products of the \(A(x_0)\).
Notice that \(\widetilde A\) is local, being a direct limit of
local rings. Thus \SGARefLink{sga3:xviii:lemma:3.12}{Lemmas 3.12} and \SGARefLink{sga3:xviii:lemma:3.10}{3.10} give the quotient \(G\)
over \(\operatorname{Spec}\widetilde A\), hence, by the preceding
remarks, over one of the \(\operatorname{Spec}A_i\), and therefore
fppf-locally.

In fact, it follows from the constructions that fppf-locally one
can find a finite set of sections
\(\{c_i\}\), \(i=1,\ldots,r\), such that \(G\) is covered by the
\(c_i\times X\). Assertions \SGARefLink{sga3:xviii:lemma:3.8:item:ii}{(ii)} and \SGARefLink{sga3:xviii:thm:3.7:item:iv}{(iv)} follow easily from this
fact, and we leave the verification of \SGARefLink{sga3:xviii:lemma:3.8:item:i}{(i)} to the reader.

\SGARefTarget{sga3:xviii:corollary:3.13}\paragraph{\SGARefLink{sga3:xviii:corollary:3.13}{Corollary 3.13}.}
\textit{The group sheaf \(G\) determined by a group germ
\((X,W)\) over \(S\) is representable in each of the following
situations:}
\begingroup
\renewcommand{\labelenumi}{\textup{(\roman{enumi})}}
\begin{enumerate}
\item \textit{\(S\) is Artinian.}
\item \textit{For every local scheme \(S'\to S\) at a closed
point \(s\) of \(S\), the scheme \(X_{S'}\) has a set of sections
whose restrictions to the closed fiber form a set with closure
of dimension \(\dim X_s\).}
\item\SGAInlineRefTarget{sga3:xviii:corollary:3-13:item:iii} \textit{\(S\) is strictly local and \(X/S\) is smooth.}
\item \textit{There exists an fppf covering \(S'\to S\) such
that \(G_{S'}\) is representable and affine over \(S\).}
\end{enumerate}
\endgroup

Indeed, \SGARefLink{sga3:xviii:lemma:3.8:item:ii}{(ii)} follows from \SGARefLink{sga3:xviii:lemma:3.10}{Lemmas 3.10} and \SGARefLink{sga3:xviii:lemma:3.12}{3.12}; \SGARefLink{sga3:xviii:thm:3.7:item:iii}{(iii)} follows
directly from \SGARefLink{sga3:xviii:lemma:3.8:item:ii}{(ii)} and Hensel's lemma; \SGARefLink{sga3:xviii:thm:3.7:item:iv}{(iv)} follows from descent
for affine schemes; and \SGARefLink{sga3:xviii:lemma:3.8:item:i}{(i)} follows from descent for group
schemes, which is possible here because every finite subset of a
group over a field is known to be contained in an affine open
subset (Exposé VI).

\SGARefTarget{sga3:xviii:bibliography}
\begin{thebibliography}{9}
\bibitem{Weil1948}
A.~Weil,
\emph{Variétés abéliennes et courbes algébriques},
Hermann, Paris, 1948.

\bibitem{Yanagihara1963}
H.~Yanagihara,
\emph{Reduction of group varieties and transformation spaces},
Journ. Sci. Hiroshima Univ., Ser. A-I, vol.~27, no.~1,
June 1963.
\end{thebibliography}
